\documentclass[12pt]{article}

% ============================================================
% PACKAGES
% ============================================================

\usepackage{amsmath}      % Mathematical environments: align, equation, etc.
\usepackage{amsthm}       % Theorem, definition, proof environments
\usepackage{amsfonts}     % Mathematical fonts such as \mathbb
\usepackage{amssymb}      % Additional mathematical symbols
\usepackage{mathtools}    % Extensions to amsmath
\usepackage{geometry}     % Page margins
\usepackage{enumitem}     % Customize enumerate/itemize
\usepackage{mdframed}     % Framed boxes
\usepackage{tikz}         % Draw diagrams
\usepackage{hyperref}     % Clickable references and links

\geometry{
    letterpaper,
    margin=0.8in
}

\linespread{1.2}


% ============================================================
% NUMBER SYSTEMS
% ============================================================
% Instead of writing \mathbb{R} every time, use \RR.

\newcommand{\RR}{\mathbb{R}}   % Real numbers
\newcommand{\CC}{\mathbb{C}}   % Complex numbers
\newcommand{\NN}{\mathbb{N}}   % Natural numbers
\newcommand{\ZZ}{\mathbb{Z}}   % Integers
\newcommand{\QQ}{\mathbb{Q}}   % Rational numbers


% ============================================================
% COMMON MATHEMATICAL NOTATION
% ============================================================
% These commands are useful for frequently used notation.

\newcommand{\set}[1]{\left\{#1\right\}}
% Example:
%   \set{x \in \RR : x > 0}
% produces:
%   { x in R : x > 0 }

\newcommand{\abs}[1]{\left|#1\right|}
% Example:
%   \abs{x-y}
% produces:
%   |x-y|

\newcommand{\norm}[1]{\left\|#1\right\|}
% Example:
%   \norm{x}
% produces:
%   ||x||

\newcommand{\inner}[2]{\left\langle #1,#2\right\rangle}
% Example:
%   \inner{x}{y}
% produces:
%   <x,y>

\newcommand{\parens}[1]{\left(#1\right)}
% Example:
%   \parens{x+y}

\newcommand{\brac}[1]{\left[#1\right]}
% Example:
%   \brac{a,b}

\newcommand{\angles}[1]{\left\langle #1\right\rangle}
% Example:
%   \angles{x,y}


% ============================================================
% THEOREM / THEORY COMPONENTS
% ============================================================

% "definition" style is upright text.
% Use this for definitions, examples, and remarks.

\theoremstyle{definition}

% ------------------------------------------------------------
% DEFINITION
% ------------------------------------------------------------
% Use when introducing a formal mathematical definition.
%
% Usage:
%
% \begin{definition}[Metric]
% A metric on a set X is a function ...
% \end{definition}
%
% [Metric] is optional and gives the definition a name.

\newtheorem{definition}{Definition}[section]


% ------------------------------------------------------------
% EXAMPLE
% ------------------------------------------------------------
% Use to give a concrete example of a definition/theorem.
%
% Usage:
%
% \begin{example}
% Consider X = \RR with the usual distance.
% \end{example}

\newtheorem{example}[definition]{Example}


% ------------------------------------------------------------
% REMARK
% ------------------------------------------------------------
% Use for observations, comments, or small additional facts.
% It is NOT intended for major mathematical results.
%
% Usage:
%
% \begin{remark}
% This definition does not require X to be a subset of \RR^n.
% \end{remark}

\newtheorem{remark}[definition]{Remark}


% ============================================================
% THEOREM-LIKE COMPONENTS
% ============================================================

% "plain" style makes the mathematical statements italic.

\theoremstyle{plain}


% ------------------------------------------------------------
% THEOREM
% ------------------------------------------------------------
% Use for an important/general mathematical result.
%
% Usage:
%
% \begin{theorem}[Cauchy--Schwarz Inequality]
% For all x,y \in \RR^n,
% \[
% |\inner{x}{y}| \leq \norm{x}\norm{y}.
% \]
% \end{theorem}

\newtheorem{theorem}[definition]{Theorem}


% ------------------------------------------------------------
% PROPOSITION
% ------------------------------------------------------------
% Use for a useful mathematical result that is usually less
% central than a theorem.
%
% Usage:
%
% \begin{proposition}
% Every finite-dimensional subspace of \RR^n is closed.
% \end{proposition}

\newtheorem{proposition}[definition]{Proposition}


% ------------------------------------------------------------
% LEMMA
% ------------------------------------------------------------
% Use for an auxiliary result whose main purpose is to help
% prove another theorem/proposition.
%
% Usage:
%
% \begin{lemma}
% ...
% \end{lemma}

\newtheorem{lemma}[definition]{Lemma}




% ------------------------------------------------------------
% COROLLARY
% ------------------------------------------------------------
% Use for a result that follows relatively directly from a
% theorem/proposition that was just established.
%
% Usage:
%
% \begin{corollary}
% ...
% \end{corollary}

\newtheorem{corollary}[definition]{Corollary}


% ============================================================
% PROOF
% ============================================================
% The proof environment comes from amsthm.
%
% Usage:
%
% \begin{proof}
% Suppose that ...
%
% Therefore, ...
% \end{proof}
%
% LaTeX automatically adds the QED symbol at the end.


% ============================================================
% CUSTOM BOX: IDEA
% ============================================================
% Use this when you want to record intuition or the main idea
% behind a mathematical concept.
%
% Usage:
%
% \begin{idea}
% \textbf{Key idea.}
% A metric gives us a notion of distance...
% \end{idea}

\newmdenv[
    linewidth=1pt,
    skipabove=10pt,
    skipbelow=10pt
]{idea}


% ============================================================
% CUSTOM BOX: WARNING
% ============================================================
% Use this for common mistakes, distinctions, or things that
% are easy to confuse.
%
% Usage:
%
% \begin{warning}
% Do not confuse continuity with uniform continuity.
% \end{warning}

\newmdenv[
    linewidth=1pt,
    skipabove=10pt,
    skipbelow=10pt
]{warning}

% ============================================================
% CUSTOM BOX: NOTE
% ============================================================
% Use this for a short personal note, useful observation,
% connection between concepts, or explanation.
%
% Unlike "Remark", this is intended for informal notes.
%
% Usage:
%
% \begin{note}
% A sequence can be viewed as a function from $\NN$ to $\RR$.
% \end{note}

\newmdenv[
    linewidth=1pt,
    skipabove=10pt,
    skipbelow=10pt
]{note}


% ============================================================
% DOCUMENT
% ============================================================

\begin{document}


% ============================================================
% TITLE
% ============================================================

\begin{center}
    {\Large \textbf{Construction of $\RR$}}\\[0.5em]
    {\large Lecture 3}
\end{center}

% Automatically generate the table of contents.
\tableofcontents

\newpage
\section{Least Upper Bound}

\begin{definition}[Upper Bound]
Given that $E \in S$ is ordered. If there is a $\beta \in S$ such that for all $x \in E$, we have $x \le \beta$, then we say that $E$ is upper bounded and $\beta$ is an upper bound of $E$.
\end{definition}

\begin{example}
Given $A = \{1,2,3\} \in \QQ$, we have that 100 is an upper bound of $A$.
\end{example}

\begin{note}
We can define lower bound similarly with upper bound.

And remember the bigger set $S$, it is important for us to know if the set $E$ has a least upper bound (will make clear in after parts).
\end{note}

\begin{definition}[Least Upper Bound]
If there is an $\alpha \in S$ satisfied:
\begin{itemize}
    \item $\alpha$ is an upper bound of $E$.
    \item If $\gamma < \alpha$, then $\gamma$ is not an upper bound of $E$.
\end{itemize}
Then $\alpha$ is the least upper bound of $E$ or the supremum of $E$. We denote as: $\sup E = \alpha$.
\end{definition}

\begin{example}
Given $S = \QQ$,

$\sup(\{\frac{1}{2}, 1, 2\}) = 2$. This example is very simple, $2$ is of course an upper bound of the set. And if there exists a $\gamma < 2$ of course $\gamma$ can't be an upper bound because $2$ is in the set.

$\sup \QQ_- = 0$. Again, very simple.

$\sup(\{r \in \QQ : r <2\}) = 2$. Trivial.

$\sup(\{r \in \QQ : r^2 <2\})$ does not exist. This example is non-trivial. You may draw it out on the real line and tell me that the $\sup$ is $\sqrt{2}$. However, we haven't learned the square root yet, and $\sqrt{2}$ is not in $\QQ$. But you may question why there is no $\sup$? The intuition is that you can always find an upper bound which is smaller than any given upper bound (just try to draw it out and imagine, it is very natural for you to think there always exist smaller upper bound).
\end{example}

\begin{definition}[Least Upper Bound Property]
A set $S$ has the Least Upper Bound Property if every non-empty subset of $S$ which has an upper bound has a least upper bound. 
\end{definition}

\begin{example}
According to the example above, $\QQ$ has been proved to be a non-"least upper bound property".
\end{example}

\section{Construction of the Real numbers}

\begin{definition}[Cut]
A cut $\alpha$ is a subset of $\QQ$, such that:
\begin{itemize}
    \item \textbf{Non-trivial} $\alpha \neq \emptyset, \QQ$.
    \item \textbf{Closed downward} If $p \in \alpha$, $q \in \QQ$, and $q<p$ then $q \in \alpha$.
    \item \textbf{No largest member} If $p \in \alpha$ then $p < r$ for some $r \in \alpha$.
\end{itemize}
\end{definition}

\begin{note}
The definition of a cut (in $\QQ$!) may confusing, but it is just an interval in $\RR$ with the form $(-\infty, a)$ with $a \in \QQ$, and we just take rational numbers. Or you can make it more mathematical: $cut = \{p\in\QQ: cond.\}$ with $cond.$ is expressed in elements in $\QQ$ (only $\QQ$!) and $cut$ has to satisfy 2 conditions above.
\end{note}

\begin{example}

\begin{itemize}
    \item $\alpha = \QQ_-$ is a cut.
    \item $\gamma = \{r \in \QQ: r \le 2\}$ is not a cut (not satisfy "No largest member")
\end{itemize}

\end{example}

\begin{definition}[The Real]
$\RR = \{\alpha\}$ such that $\alpha$ is a cut. We need to equip $\RR$ with arithmetic, order, in a way that $\RR$ extends $\QQ$.
\end{definition}

\begin{definition}[Order in $\RR$]
$\alpha < \beta$ if $\alpha \subset \beta$ (proper subset).
\end{definition}

\begin{note}
Just draw some sets out (intervals on line, not exactly), you can confirm (by intuition) that the order defined above satisfied trichotomy and transitive of an order.     
\end{note}

\begin{definition}[Addition in $\RR$]
$\alpha + \beta = \{r+s: r\in\alpha, s\in\beta\}$
\end{definition}


\begin{note}
We have to check if the defined addition satisfies 5 axioms of addition in field. Again, remember, conditions for elements in a cut have to write in form of elements in $\QQ$.

With closeness, commutative, associative, it is easy, so I'll skip.

With additive identity, we have $\mathbf{0}^* = \{p\in\QQ: p<0\}$ is the additive identity. Given any $\alpha$, we have $\alpha + \mathbf{0}^* = \alpha$. Intuitively, when adding $\mathbf{0}^*$, we always decrease the value of a given element, then it still in the old cut. And if we want to recreate any element $a \in \alpha$, we just need to take $a + \epsilon \in \alpha$ (no largest member) with positive, large enough $\epsilon$, then add with $-\epsilon$ to return $a$.

With additive inverse, given any $\alpha = \{p \in \QQ: p < r\}$, we have $-\alpha = \{p \in \QQ: \exists r \notin \alpha, p < -r\}$. \textbf{Please remember to draw it out.}
\end{note}

\begin{definition}[Multiplication in $\RR$]
We gonna define multiplication for positive cuts first $\{\alpha: \alpha >\mathbf{0}^*\} = \RR_+$

Let $\alpha, \beta \in \RR_+$, then the multiplication $\alpha \cdot \beta = \{p \in \QQ: p < rs, \exists r \in \alpha, s\in\beta, r>0, s>0\}$.
\end{definition}

\begin{note}
The multiplication of positive cuts in $\RR$ is just multiply all the positive part (ignore the 0 and negative part), then keep the cut structure intact.

Given multiplication in positive, we can extend to any sign:
\begin{itemize}
    \item $\alpha = \mathbf{0}^*$ or $\mathbf{0}^* = \beta$, then $\alpha\beta=0^*$.
    \item $\alpha, \beta>\mathbf{0}^*$, then use the definition above.
    \item $\alpha > \mathbf{0}^*$, $\beta < \mathbf{0}^*$ then, $\alpha\beta = -(\alpha(-\beta))$.
    \item $\alpha < \mathbf{0}^*$, $\beta > \mathbf{0}^*$ then, $\alpha\beta = -(-\alpha(\beta))$.
    \item $\alpha < \mathbf{0}^*$, $\beta < \mathbf{0}^*$ then, $\alpha\beta = -\alpha(-\beta)$.
\end{itemize}
\end{note}

\begin{note}
Again checking the multiplication axioms, I'll focus on multiplicative identity and multiplicative inverse.

About multiplicative identity, it is $\mathbf{1}^* = \{p\in\QQ:p<1\}$. Now, we need to show $\mathbf{1}^* \cdot \alpha = \alpha$. The key thing in the whole multiplication stuffs is the intuition of multiply positive part while preserving the cut's properties. Then you want to ask, how multiplying the positive part of a intervals (not exactly) to $(0,1)$ (not exactly), can keep the positive part intact? Consider the value of the new "interval", of course the value stay the same, because they multiply to $(0,1)$. Now, if you want to recreate any element $x$ in the positive part, simply take $x + \epsilon$ (large enough, positive $\epsilon$ - no largest member), then multiply it with $\frac{x}{x+\epsilon} \in \mathbf{1}^*$, then you obtain $x$. Hence, you can keep the positive part intact.

About multiplicative inverse, very similar to additive inverse, we have $\alpha^{-1} = \{p \in \QQ: \exists r \notin \alpha, r>0, p < \frac{1}{r}\}$. The intuition of this, I think the best way is to \textbf{draw it out}.
\end{note}

\begin{note}
And last, but not least, distributive axiom. We have to check that. Oh forget it, accept that, I can't imagine.
\end{note}


\begin{theorem}
$\RR$ is an ordered field with the Least Upper Bound Property, and $\RR$ contains $\QQ$ as a subfield.
\end{theorem}

\end{document}